% Nouveau EGA V (cote 25): archive cote 25, site page 5 (PDF page 6).
% typescript. Draft compared: none.
% AI-transcribed by Claude Opus 5.5 (claude-opus-5-5, adaptive thinking, effort high), Message Batches API, batch ega25b (msgbatch_01A7n5YuXubUPNK1zsepi8RG), 2026-11-13; not yet checked by eye.
\scanpage{25/5}
\mspage{2}

ou encore $\underline{L}^{-1}\subset\underline{E}$) et $\emptyset$ s'obtient en prenant\inserted{pour $\emptyset$} une section de $\underline{L}^{-1}$ qui
ne s'annule en aucun point \xout{\illegible{1}}, i.e. une trivialisation de $\underline{L}^{-1}$
(qui existe en tous cas localement). Notons d'ailleurs que $H_{\hand{\xi}}$ n'est
autre que $P(\underline{E}/\underline{L}^{-1})$ (isomorphisme canonique), ce qui est une troisième
façon de décrire $H_{\hand{\xi}}$ (NB $P(\underline{E}/\underline{L}^{-1})$ est bien plongé canoniquement dans
$P = P(\underline{E})$ ), qui a l'avantage de montrer de plus que $H$ est un fibré projec-%
tif sur $S$, et à fortiori est lisse sur $S$ (encore aurait-il fallu dire \xout{\unsure{qu}}
\margin{$\|$} au par. 17 qu'un fibré projectif est lisse ...). Il serait sans doute
\margin{$\|$} préférable de commencer par celle-là.
\uline{Remarques}. La formation de $H_{\hand{\xi}}$ en termes de $\hand{\xi}$ est compatible avec
le changement de base, comme on constate aussitôt., en d'autres termes,
on trouve un homomorphisme de foncteurs $(\mathrm{Sch})^{\circ}_{/S}\to(\mathrm{Ens})$
\[ \check{P} \longrightarrow \underline{\mathrm{Div}}\,(P/S) \qquad , \]
où le deuxième membre désigne le foncteur des ``diviseurs relatifs'' de $P/S$,
dont la valeur en $S'$ ($S$-préschéma arbitraire) est l'ensemble des
sous-schémas fermés de $P_{S'}$ qui sont intersection complète transversale
de codimension 1 relativement à $S'$ (cf par.19). Il est facile de montrer
que cet\xout{\illegible{1}} homomorphisme de foncteurs est un monomorphisme, en d'autres
termes que $\hand{\xi}$ est déterminé par $H_{\hand{\xi}}$ . (Ce dernier fait justifie la termino-%
logie ``schéma des hyperplans'' utilisée plus \xout{\unsure{a}}haut). Nous verrons \xout{\illegible{1}} \margin{)}
\xout{\unsure{Chap.V}} que le foncteur $\underline{\mathrm{Div}}$ $(P/S)$ est représentable par \xout{\unsure{un}}\inserted{le pré}schéma somme
des $P(\mathrm{Symm}^{k}(\check{E}))$, de sorte que $\overset{\hand{\vee}}{P}$ s'identifie à un sous-préschéma ouvert
et fermé de $\underline{\mathrm{Div}}$ $(P/S)$ ... (NB \over{A}{À} vrai dire, la d\over{e}{é}termination des diviseurs
\margin{$|$} relatifs\inserted{de $P/S$} pourrait se faire avec les moyens du bord dès maintenant, en
\margin{$|$} utilisant les résultats du Chap III, et pourrait être ajoutée \xout{\illegible{1}} à titre
\margin{$|$} d'exemple au par.19 ...).

\mspage{5}
